\section{Identification with zero support: complete proofs}
\label{sec:annex-p16-identification-with-zero-support}

\begin{annexitem}{annex:p16.zero-support-quantifiers}{thm:p16-zero-support-quantifiers}{Joint-action quantifiers and the exact positivity implication}
\begin{proof}
For \(I=N\), its outside parent set is empty. The empty tuple
has probability one, so either support condition for a full action
assignment \(v\) says exactly \(p(v)>0\). Requiring that
condition for every full assignment therefore implies strict full
positivity. This argument uses no Markov or independence assumption.

Conversely suppose every full mass is positive. For any action set
\(I\), its outside parent set \(B_I\) is disjoint from it.
Every proposed \((i,b)\) extends to a full assignment because
all value sets are nonempty. The marginal \(p(i,b)\) contains
that positive full mass, so it is positive. Consequently both support
conditions hold for every \(i\) and every \(I\). This proves
both equivalences, including the empty observed graph.

For the stated two-variable example the only positive masses are
\(p(0,0)=p(1,1)=1/2\). Both observed parent sets are empty.
Every singleton action value therefore has conditional probability
\(1/2\) on its sole empty parent stratum, satisfying both readings.
The two off-diagonal full masses are nevertheless zero. In this
two-vertex graph every proper action set is a singleton or empty,
and the empty action also satisfies both conditions. Thus even all
proper action sets do not force full positivity. A single
original action set can be one of these singletons, proving that
its support conditions alone also do not force full positivity.
In general \(\mathcal A_{\mathrm{all}}\subseteq
\mathcal A_{\mathrm{obs}}\). A member of the latter has positive
unconditional action mass, because at least one parent stratum has
positive mass and its joint mass with that action must be positive.
None of these observations restricts the existence of an equation-
replacement intervention at another value.
\end{proof}
\end{annexitem}

\begin{annexitem}{annex:p16.zero-front-door-obstruction}{thm:p16-zero-front-door-obstruction}{The missing mediator strata and two distinct intervention laws}
\begin{proof}
The shared root \(U\) is declared at \(A,Y\); the root
\(E\) is private at \(Y\). The equation at \(M\) uses
only its observed parent \(A\). In model 2 the \(Y\) equation
ignores its declared \(M,U\) arguments, as the graph-defined
class permits. Thus both models have the required declared graph
and independent finite latent inputs.

Observationally \(A=M=U\), so \(M\oplus U=0\) and
both models have \(Y=E\). For each \(a\), the root value
\(U=a\) has probability \(1/2\), and the response probabilities
are \(3/4,1/4\). This gives \(3/8,1/8\) on the two
diagonal response cells and zero when \(m\ne a\). Their
total mass is one. The original action \(A\) has no observed
parent, and \(p(A=a)=1/2\) for both values. Both support
readings therefore hold for every original action value.

Replacing \(A\) by \(a\) also makes \(M=a\), while
\(U\) retains its fair prior law. The bit \(a\oplus U\)
is fair and independent of \(E\). Thus model 1 has response-one
probability \(1/2+(1/2)(1/4)=5/8\); model 2 retains
probability \(1/4\). These two compatible models prove
nonidentification at exactly the stated observed law and supported
original actions.

For completeness, no original-query hedge is present in this graph.
A bidirected-connected forest containing the action vertex \(A\)
can use only \(A\) or \(\{A,Y\}\), since \(M\) has
no incident bidirected edge. A singleton \(A\) has no smaller
nonaction forest with the same roots. The induced directed graph
on \(A,Y\) has no arrow; any selected forest there has both
vertices as roots, so a smaller forest excluding \(A\) cannot
have the same root set. These exhaust the possibilities.

The ID graph trace splits the post-action districts \(\{M\}\)
and \(\{Y\}\). The former yields \(p(m\mid a)\).
For the latter it extracts the full-graph district \(\{A,Y\}\)
and then sums \(A\), giving the usual formal expression
\[
 \sum_m p(m\mid a)\sum_{a'}p(a')\theta_Y(1\mid a',m).
\]
Tian's graph procedure makes the same two district computations.
Here \(p(m\mid a)\) is one at \(m=a\) and zero otherwise.
At the supported pair \((a,a)\), the last conditional is
\(1/4\). At the pair \((1-a,a)\) its denominator is zero,
so a normalized binary version has an arbitrary response-one value
\(t\in[0,1]\). The expression becomes \(1/8+t/2\).
Its terms at \(m\ne a\) may indeed be omitted because their
outer factor is zero. The term involving \(a'=1-a,m=a\)
has the positive multiplier \(p(a')=1/2\), however. Its
zero observed joint mass cannot replace that different multiplier.
For the complementary action \(M=a\), the parent tuple
\(A=1-a\) also directly fails both intermediate support readings.
Thus original-action support does not license this required extraction.
\end{proof}
\end{annexitem}

\begin{annexitem}{annex:p16.zero-version-factorization}{thm:p16-zero-version-factorization}{Every prefix case and a simultaneous structural version}
\begin{proof}
For each fixed prefix, summing \(\rho_j\) over its next
coordinate gives \(\rho_{j-1}\). If the latter is positive,
the ratio is nonnegative and normalized. If it is zero, every
nonnegative summand \(\rho_j\) is zero, so any normalized
next-coordinate table satisfies the displayed equality. Choosing
such tables for each of the finitely many rows constructs versions.

Fix a full assignment. If all its prefixes are positive, the
product of ratios telescopes to \(q_R\). Otherwise its first
zero prefix \(j_0\) has a positive predecessor and therefore
\(\theta_{j_0}=0\). The later version values are finite
numbers in \([0,1]\), so the whole product is zero, exactly
\(q_R\) at that assignment. Grouping this finite product by
district proves the first identity even at zero cells. To normalize
\(h_D^\theta\) on \(D\), sum its last topological coordinate
first. Earlier factors do not use that later coordinate, and its
own conditional table sums to one. Repeating in reverse order
leaves one. The outside coordinates are fixed throughout.

We now construct versions matching a given structural model. Keep
the maximal districts of the whole current set \(R\), and
partition its independent latent roots into their actual district
groups \(\Lambda_D\). Define
\[
 g_{D,j}(v)=\sum_{u_{\Lambda_D}}
       \prod_{\alpha\in\Lambda_D}r_\alpha(u_\alpha)
       \prod_{i\in D\cap R_j}
       \mathbf1_{\{v_i=f_i(v_{\operatorname{pa}(i)},u_{\Lambda_i})\}}.
\]
All fixed outside parent values are retained in these equations.
The latent-sum foundation proves
\(\rho_j=\prod_D g_{D,j}\), \(g_{D,0}=1\),
\(g_{D,k}=Q[D]\), and that only the district containing
\(r_j\) changes at step \(j\). These identities follow
by finite distributivity and require no positivity.

For \(r_j\in D\), sum \(g_{D,j}\) over \(v_{r_j}\).
All earlier indicators are independent of this later observed
coordinate. The new total equation has exactly one permitted
output for each fixed parent tuple and latent tuple, so its
indicator sums to one. Hence that sum equals \(g_{D,j-1}\).
When \(g_{D,j-1}>0\), choose
\(\theta^M_j=g_{D,j}/g_{D,j-1}\). When it is zero,
all corresponding \(g_{D,j}\) values are zero and choose any
normalized row. These choices are functions of the observed prefix
and fixed outside values, so they are valid conditional tables.
If \(\rho_{j-1}>0\), all its district factors are positive,
and cancelling the unchanged ones shows
\(\theta^M_j=\rho_j/\rho_{j-1}\). If \(\rho_{j-1}=0\),
both sides of the defining prefix identity are zero. Thus
\(\theta^M\) is a conditional version of \(q_R\).

Within each district the product of these chosen factors equals
\(g_{D,k}\). If its successive district factors stay positive,
they telescope from one. If a first factor becomes zero, its
predecessor is positive and the corresponding chosen conditional
is zero; all later district factors remain zero at this full
assignment, since an additional indicator cannot increase their
nonnegative latent sum. The product is then zero on both sides.
This proves \(h_D^{\theta^M}=Q[D]\) for every district and
assignment simultaneously, including all zero cases.

Finally consider arbitrary versions at the fixed assignment in
the statement. Up to and including the first zero prefix, the
relevant conditionals are fixed ratios. After that point every
next row has a zero predecessor and is free, subject only to
normalization. A singleton row is necessarily one. If the product
of the fixed district factors is zero, every completion gives
zero. If that fixed product is positive and there are no later
non-singleton district coordinates, the product is fixed. Otherwise
choose every free district row on this assignment to concentrate
on its assigned coordinate, giving the positive fixed product.
In another completion choose one non-singleton such row to
concentrate on a different value, giving zero. Rows at different
indices are separate choices, so both completions are legitimate.
The product is then not version-invariant. This proves the finite
decision rule and, through the model-compatible version, every
claimed unconditional recovery of a district cell.
\end{proof}
\end{annexitem}

\begin{annexitem}{annex:p16.zero-masked-extraction}{thm:p16-zero-masked-extraction}{The complementary-action factor, the mask and the finite table operation}
\begin{proof}
Fix the current outside parameters and the requested assignment
\(j\) on \(J=R\setminus D\). Because \(D\) is a
maximal district, the latent-input product theorem gives
\[
 q_R(v_D,j)=Q[D](v_D;j)\,Q[J](j;v_D).
\]
Indeed the other maximal district factors multiply to \(Q[J]\),
by the same theorem on \(J\). The kernel \(Q[J]\) uses
outside observed values in \(D\) only at the observed parents
\(B=\operatorname{pa}_{G[R]}(J)\setminus J\). Therefore,
for each \(b\in S_B\), summing the preceding identity gives
\[
 q_R(j,b)=Q[J](j;b)\sum_{v_{D\setminus B}}Q[D](v_D;j).
\]
The latter sum is nonnegative and at most one because \(Q[D]\)
is normalized on all of \(D\) at \(j\). Under the overlap
hypothesis the left side is positive whenever \(q_R(b)>0\),
so \(Q[J](j;b)>0\) on every observed parent stratum.

We must also control parent tuples absent from the observed law.
Suppose \(Q[D](d;j)>0\). There is a positive-prior tuple
\(u_D\) of the district's latent inputs producing \(D=d\)
when \(J=j\). Fix this tuple throughout the following argument.
Because \(D\) is maximal, its latent inputs are disjoint from
those used by \(J\), and the two groups are independent. Choose
any positive-prior tuple \(u_J\) for the latter group and solve
the current model naturally. The resulting full assignment has
positive observed mass, so its parent subtuple \(b\) satisfies
\(q_R(b)>0\). Overlap gives \(q_R(j,b)>0\), hence
\(Q[J](j;b)>0\). There is therefore a positive-prior tuple
\(u'_J\) satisfying all equations on \(J\) at the fixed
values \(j\) and external parents \(b\).

Order the vertices of \(J\) by the restricted topological order.
If their current natural values already equal \(j\), stop.
Otherwise let the first differing vertex be the next one to repair.
Resolve the full model using the same \(u_D\) and the new
\(u'_J\). All coordinates before that first difference remain
unchanged: at vertices in \(D\), the latent tuple and earlier
observed parents are unchanged; at earlier vertices in \(J\),
the new tuple satisfies their equations at their target values,
which already equal their old values, and at the same earlier
\(D\) parents. This is induction in the full topological order.
At the first differing \(J\) vertex, the same argument makes
its new value the required one. Thus the initial matching segment
of \(J\) strictly lengthens. The new natural assignment still
has positive mass, by independence and positivity of the two latent
tuples, so its possibly changed parent stratum again satisfies
the overlap hypothesis.

Repeating repairs reaches a natural solution with \(J=j\) in
at most \(|J|\) repairs. Its \(D\) solution must be \(d\),
by uniqueness with the fixed \(u_D,j\). Hence
\(q_R(d,j)>0\). Inputs with no current child may be given any
positive-prior values and do not affect this argument. This proves
that every positive district cell on the intervention slice has
positive full observed mass there. Its contrapositive proves the
required zero case, including formerly null parent tuples. When
\(J\) is empty the same conclusion is immediate from \(q_R=Q[D]\).

At a positive full \(q_R(v_D,j)\), every prefix is positive,
and the already proved district recovery formula identifies
\(Q[D]\) as the product of its ordinary conditional ratios.
This proves both cases of the masked formula. No ratio is evaluated
in its zero case, so arbitrary versions there cannot affect it.
The case \(J=\varnothing\) uses \(Q[\varnothing]=1\)
and the same argument.

For the more general recovery operation, every cell determined
by version invariance equals the actual district cell in every
compatible model by the simultaneous structural version theorem.
If cells remain undetermined, they have zero full observed mass;
otherwise all their prefixes would be positive. The overlap
argument just given fills those cells by their forced actual zero
when its hypothesis holds. Thus every value written into a returned
table equals the actual structural kernel value. If the whole
table is determined, it inherits normalization and nonnegativity
from that kernel. If some cells remain, the operation reports
them and does not pass the incomplete table to recursion.

All steps are finite: prefix sums, positivity comparisons, a scan
for the first zero, inspection of finite value-set cardinalities,
and, if needed, the finitely many parent-stratum sums in the
overlap test. Division occurs only at explicitly positive prefixes.
In particular \(q_R(j,b)>0\) implies \(q_R(b)>0\), so
the overlap test is exactly the requirement of a defined positive
conditional for this complementary action at every observed parent
stratum in the current slice. The repair argument proves why no
assumption on null parent tuples is needed. It does not assert positivity of
every full \(q_R\) cell or of every unrelated intervention value.
\end{proof}
\end{annexitem}

\begin{annexitem}{annex:p16.zero-certain-action}{thm:p16-zero-certain-action}{Pointwise consistency on an event of probability one}
\begin{proof}
Use the restricted structural model whose law at the fixed outside
parameters is \(q_R\). On every latent outcome for which its
natural \(X\) coordinates equal \(x\), the natural solution
also satisfies all equations after replacing the \(X\) equations
by \(x\). Uniqueness of the acyclic solution makes the two
solutions identical on all current coordinates at that outcome.
The event has probability one by hypothesis. Consequently every
response event has the same probability before and after replacement,
which gives \(P^R_x(y)=q_R(y)\). Outcomes of zero latent
mass do not affect this equality. No independence between the
action and response variables, or between their shared latent
inputs, was assumed in this consistency argument.
\end{proof}
\end{annexitem}

\begin{annexitem}{annex:p16.zero-algorithm-soundness}{thm:p16-zero-algorithm-soundness}{Every scalar branch, its support, and each discarded zero term}
\begin{proof}
Each request carries a full known current table at one fixed
outside parameter tuple. Initially this is the observed table.
An ancestral marginal remains a full normalized nonnegative table
and equals its structural \(Q\) kernel by the unconditional
latent-sum theorem. Every table returned by \(\mathsf{Recover}\)
is a full actual district kernel on its stated slice by the
recovery theorem. These are the only new tables passed to recursion.
Thus induction establishes the actual-kernel invariant at every
reached request; no conditional observation of an outside tuple
has replaced its intervention meaning.

For ID use its measure \(|R|+|R\setminus X|\). Every
recursive graph edge still strictly decreases it, as already
proved: ancestry or larger-district restriction decreases the
random set, extra actions decrease the nonaction set, and a split
child has a proper smaller nonaction district. Trying different
reference assignments gives finitely many children of the same
smaller measure. Summation loops and all recovery operations are
finite. The certain-action shortcut and unresolved returns terminate
immediately. Hence the scalar ID procedure terminates. For each
component procedure \(|T|-|C|\) strictly decreases whenever
it recurses, and the outer Tian procedure has finitely many terms,
factors and reference choices. It also terminates.

We prove soundness by induction on these measures. The certain-action
return is the preceding consistency theorem. ID's empty-action
case is ordinary marginalization, included in that return. At its
ancestral branch the known marginal is the correct smaller-model
law, and the original and smaller intervention targets agree by
the unconditional ancestral reduction proof. Any newly displayed
outside coordinates can be set arbitrarily when they are not
original action coordinates, since that proof establishes their
irrelevance. An inductively known child therefore gives the right
parent answer.

At ID's additional-action branch, the already proved structural
nonancestor identity gives \(P^R_x(y)=P^R_{x,w}(y)\) for
every reference tuple \(w\), without a conditional probability
division. Each known child answer is therefore correct for the
same parent target. Answers from different successful choices
must agree because they equal that target in every compatible
model. If none is known the retained support records assert no value.

At the split branch the unconditional district product identity
expresses the target as its specified finite sum of products of
actual district masses. An inductively known child supplies the
correct factor. If one factor of a given term is known to be
zero, the true term is zero in every compatible model because
all other factors are finite probabilities. Those other factors
need not be computed and their unresolved records are not numerical
operands. If no factor is proved zero, the rule returns a product
only when all factors are known; the product is then correct.
The parent returns a sum only when every term is determined in
one of these two ways. Thus every returned sum equals the true
target. A zero in the input observed table is not by itself a
zero in one of these required causal factors.

At ID's graph-failure branch an unresolved graph record makes
no numerical claim. At its maximal-district base, recovery first
provides the full actual kernel on the required complementary
action slice, or the request stays unresolved. A known return
is its correct response marginal. At the larger-district branch,
the same recovery theorem supplies the full smaller input kernel.
The original action values on the removed vertices are retained.
Combination of those fixed actions with the remaining actions
is the original intervention, so induction proves the returned
child value correct for the parent. These cover all seven ID
branches and the added consistency shortcut.

For c-identify, the ancestral base is the exact ancestral sum
of its known kernel. A graph failure remains unresolved. In the
recursive branch the full ancestral marginal is known and correct;
recovery either supplies the actual full \(Q[T']\) slice or
reports which cells are missing. Only in the former case is a
child called, so induction proves any known return equal to
the requested cell of \(Q[C]\). The optional certain-action
shortcut applies to the current complementary action and has
the same consistency proof as before.

The outer Tian procedure first recovers each required containing
district at the actual outside values for that scalar factor.
When recovery succeeds, the component argument just given proves
its known factor correct. The structural target identity is
\(P_x(y)=\Sigma_{D\setminus Y}\prod_i Q[D_i]\), valid
at zero masses by finite latent sums. The same product and sum
reasoning therefore proves every known outer answer. Values
outside \(D\cup X\) do not enter the target, by the proved
ancestor reduction and parent-dependence property of \(Q[D]\).
Any successful reference choice consequently gives the same target.

Every numerical division used by these procedures occurs inside
recovery at a prefix explicitly proved positive. Zeros filled
by overlap or version invariance have their own proof and are
never defined as \(0/0\). Positive full current cells also
license the distinct source conditional-factor derivation: dropping
equation indicators shows that every aligned partial observation
in a further-intervened model has mass at least that cell's
positive mass. Thus the source Lemma 4 rule applications on
these terms have their required support. Its unconditional
ancestral and additional-action identities remain available with
empty conditioning sets, whose mass is one; its larger-district
step combines fixed actions exactly as before. Zero cases are
settled by the separate latent-product and recovery arguments,
not by equating arbitrary null conditional versions.

Finally, if all requested response masses are known, each equals
the actual mass in every compatible model. Their vector is therefore
the actual normalized response law and is identified at the input
table. This is a soundness conclusion for known returns, without
classifying an unresolved result as a nonidentification theorem.
\end{proof}
\end{annexitem}

\begin{annexitem}{annex:p16.zero-mass-do-rules}{thm:p16-zero-mass-do-rules}{Cross-multiplication with every null-stratum case retained}
\begin{proof}
Use the full explicit-root DAG so that the already proved finite
do-calculus theorems apply, with their exact graph premises and
the stipulated assignment tuples. For Rule 1, if \(p_x(z,w)>0\),
then \(p_x(w)>0\) and the conditional equality in
\cref{thm:p16-do-rule-one} gives the first mass identity after
multiplication by these two positive denominators. If
\(p_x(z,w)=0\), then \(p_x(y,z,w)=0\) by nonnegativity,
and both sides of the mass identity are zero.

For Rule 2, if \(p_x(z,w)>0\), the theorem
\cref{thm:p16-do-rule-two} also supplies \(p_{x,z}(w)>0\).
Multiplication of its equality by the two denominators gives
the second identity. If \(p_x(z,w)=0\), then again
\(p_x(y,z,w)=0\); both sides are zero regardless of the
intervened conditioning mass. These cases are exhaustive.

For Rule 3, if \(p_x(w)\) and \(p_{x,z}(w)\) are both
positive, \cref{thm:p16-do-rule-three} gives the third identity
by multiplication. If either is zero, its corresponding joint
mass with \(y\) is zero, making both products zero. The
graph premise continues to use exactly
\(Z\setminus\operatorname{an}_{\widehat G_{\overline X}}(W)\);
no incoming arrows at the other proposed action vertices are
silently removed. All three arguments include empty tuples and
conditioning sets.

When a multiplier is zero, the resulting equality can be
\(0=0\) and gives no equality between independently assigned
conditional versions on that null event. Only multiplication
has been used in the null cases. Dividing there would not
be a permitted reversal of this proof. This establishes the
precise support boundary of the preserved do-calculus route.
\end{proof}
\end{annexitem}

\begin{annexitem}{annex:p16.zero-licensed-completeness}{thm:p16-zero-licensed-completeness}{The licensed-domain theorem and the disproved larger implication}
\begin{proof}
The domain definition is finite for each specified alphabet family
and table. There are finitely many original action assignments
and response assignments. Each scalar graph call tree is finite
by the existing decreasing measures; a graph failure is a terminal
record. Expand every split child for this domain check even if
another child gives such a record. All requested input kernels
come from ancestral sums and recovered district tables, never
from a child's returned scalar expression, so this finite expansion
defines its table requests without needing an unknown child's value.
If a recovery request fails, the law is outside the licensed
domain; otherwise all its computed tables are actual kernels.

Every strictly positive compatible law belongs to the domain.
The previous positivity invariant makes every reached current
kernel positive at every relevant assignment, so its prefixes
and every district product are fixed; recovery succeeds. This
also proves that the domain contains each positive four-valued
hedge witness. The zero-outcome example below shows that the
domain need not require full observed positivity.

Suppose the complete original query has no hedge. The proved
graph-uniform result makes either chosen ordinary graph procedure
succeed. On a licensed law every extraction table required by
that graph trace is recovered. Induct along the finite graph
computation: its base values, child returns, products and sums
are then all known. The support-aware soundness theorem identifies
them with the requested effect for every original supported
action and response value. Thus, for every finite alphabet
family, every compatible licensed law, every two compatible models
sharing that law, and every \(x\in\mathcal A_{\mathrm{obs}}(p,X)\),
the response laws agree. This is the asserted uniform positive
direction on the stated nonpositive domain.

Conversely suppose a hedge exists. The already proved construction
gives two positive four-valued models on the same graph with the
same observed law and different original intervention kernels.
Their common law is licensed by the preceding paragraph. Every
action value satisfies both original support readings because
the full law is positive. These models violate one instance of
the outer finite-alphabet/law quantifier and therefore disprove
uniform identification on the licensed domain. This proves the
negative direction and, by the earlier graph equivalences, the
stated success/failure criterion for either chosen procedure.

On the larger input-only support domain, the front-door theorem
provides two models whose original action satisfies even the
all-parent-assignment reading, yet whose effects differ in a
graph with no hedge. Its returned graph formula needs an
unobserved mediator stratum with nonzero target weight. This
is an explicit counterexample to the larger sufficient condition,
and also disproves the implication from original-action support
to the licensed-trace requirements. It is not a claim that an
intervention is structurally undefined or that the positive-law
graph theorem fails.

If instead support is required for all values of every joint
action, the first theorem of this section forces full positivity
and the earlier complete positive-law theorem applies. These
are different quantified premises. A support condition imposed
on each complementary action in its current model at every
observed positive parent stratum gives the overlap test directly, and hence
licenses those extractions; it cannot be inferred just from
the original action's condition. The proof therefore retains
each interpretation with its actual scope.
\end{proof}
\end{annexitem}

\begin{annexitem}{annex:p16.zero-worked-identification}{thm:p16-zero-worked-identification}{Every supported cell and every zero-product decision in the examples}
\begin{proof}
For the first model, the natural solution satisfies \(B=A\).
At \(A=0\), \(Y=E\), and at \(A=1\), \(Y=1\oplus E\).
Multiplying the fair action mass by the two noise masses gives
the four displayed positive cells; all assignments with \(b\ne a\)
are impossible. The total is one.

For \(P_a(Y=y)\), ID retains all directed ancestors, adds
no nonancestor action, and splits the post-action graph's two
singleton districts \(\{B\}\) and \(\{Y\}\). Its B
child first restricts to \(\{A,B\}\). At the fixed action
\(A=a\), the prefix \(p(A=a)=1/2\) is positive. The
B conditional is one at \(b=a\) and zero at \(b\ne a\).
At the latter cell the first zero is created at the B index
itself, so the version-invariance rule proves the zero. Recovery
therefore supplies the full deterministic B kernel on that slice.

For the Y factor at \((a,b)\), use order \(A,B,Y\).
If \(b=a\), its preceding prefix has mass \(1/2\), and
the two Y probabilities are the known noise probabilities.
If \(b\ne a\), the prefix first becomes zero at B, outside
the singleton Y district. Its Y version is free and the overlap
test has empty parent set but joint action mass
\(p(A=a,B=b)=0\). Recovery correctly reports the Y table
on this slice unresolved. In the outer sum its B multiplier
has already been proved zero, so this whole term is zero in
every compatible model and needs no Y version. The single
remaining term \(b=a\) has B factor one and the known Y
factor. Thus \(P_0(Y=1)=1/4\) and \(P_1(Y=1)=3/4\).

Tian's outer procedure has the same post-action districts. Its
full-graph districts are singletons. Recovery of the B table at
its outside parameters again succeeds by the same first-zero
rule, even if its overlap test fails at the other B value.
Recovery of the Y table has exactly the two cases just calculated.
The component calls on recovered singleton districts return their
tables immediately. Its outer sum therefore makes the same
legitimate zero-product deletion and gives the same result.
This shows why a useful scalar evaluation can succeed even
outside the stricter domain requiring every unused table recovered.

For the second model, \(p(A=a,Y=0)=1/2\) and
\(p(A=a,Y=1)=0\) for either action value. The Y extraction
has complementary set \(\{A\}\), no outside observed
parents of that set, and overlap mass \(p(A=a)=1/2>0\).
At Y=0 the ratio is one. At Y=1 the first zero occurs at Y,
so both the version rule and the overlap mask give zero without
division by a zero prefix. The full table is recovered at each
action value in both graph procedures; there are no other
unsupported table requests. Hence this nonpositive law lies in
the licensed-trace domain, and the displayed deterministic effect
is identified by the complete known-return proof.

For the final bow example, the observed action A=0 has probability
one. The certain-action theorem gives \(P_0(Y=y)=p(Y=y)=1/2\)
for each binary response, in every compatible model with this
observed law. ID's shortcut returns these masses. Tian first
extracts its only full-graph district, which is the whole observed
table; its component shortcut then makes the same return. Without
that shortcut the graph trace has \(Y\) as its nonaction
district and \(A,Y\) as one full district, so ID returns
the raw bow failure and c-identify reaches its analogous failure.
This does not refute graph-uniform nonidentification of the bow:
the specified sparse law and supported action are a narrower
question, settled here by probability-one consistency.
\end{proof}
\end{annexitem}
